\documentclass[handout_subfiles_root.tex]{subfiles}
\usepackage[rm2]{lecturenote}

\begin{document}
\HandoutTitle{8}{The Discrete Fourier Transform (DFT)}{September 17, 2026}
\maketitle

\HandoutMarkSection{1}{The Discrete Fourier Transform (DFT)}

\begin{defbox}[title={Def}]
The DFT is a mapping of a length-$N$ sequence $x[n]$, $0\le n\le N-1$, into a length-$N$ Fourier representation $X[k]$, $k=0,\dots,N-1$, according to
\begin{equation*}
X[k] = \sum_{n=0}^{N-1} x[n]\,e^{-j\frac{2\pi kn}{N}}.
\end{equation*}
\end{defbox}

\noindent It is a \textbf{sampled version of the DTFT} for a finite-length sequence:
\begin{equation*}
X[k] = X(e^{j\omega})\Big|_{\omega=\frac{2\pi k}{N}}.
\end{equation*}
$\Rightarrow$ \HandoutHighlightYellow{discrete in both time and frequency}.

\begin{defbox}[title={Inverse DFT}]
\begin{equation*}
x[n] = \frac{1}{N}\sum_{k=0}^{N-1} X[k]\,e^{+j\frac{2\pi kn}{N}}, \qquad n=0,\dots,N-1.
\end{equation*}
\end{defbox}

\HandoutMarkSubsection{2}{Twiddle-Factor Notation}

With $W_N = e^{-j\frac{2\pi}{N}}$,
\begin{align*}
X[k] &= \sum_{n=0}^{N-1} x[n]\,W_N^{kn}, \qquad 0\le k\le N-1, \\
x[n] &= \frac{1}{N}\sum_{k=0}^{N-1} X[k]\,W_N^{-kn}, \qquad 0\le n\le N-1.
\end{align*}

\paragraph{Verification of the inverse formula.}
\begin{align*}
X[\ell] &= \sum_{n=0}^{N-1} x[n]\,W_N^{\ell n}
= \sum_{n=0}^{N-1}\left[\frac{1}{N}\sum_{k=0}^{N-1} X[k]\,W_N^{-kn}\right] W_N^{\ell n} \\
&= \sum_{k=0}^{N-1} X[k]\,\underbrace{\left[\frac{1}{N}\sum_{n=0}^{N-1} W_N^{-(k-\ell)n}\right]}_{\HandoutUnderbraceLabel{(\ast)}}.
\end{align*}
Summing the geometric series,
\begin{align*}
(\ast) &= \frac{1}{N}\,\frac{1-W_N^{-(k-\ell)N}}{1-W_N^{-(k-\ell)}}
= \frac{1}{N}\,\frac{1-e^{j2\pi(k-\ell)}}{1-e^{j\frac{2\pi(k-\ell)}{N}}} \\
&= \frac{1}{N}\,\frac{e^{j\pi(k-\ell)}\left(e^{-j\pi(k-\ell)}-e^{j\pi(k-\ell)}\right)}{e^{j\frac{\pi(k-\ell)}{N}}\left(e^{-j\frac{\pi(k-\ell)}{N}}-e^{j\frac{\pi(k-\ell)}{N}}\right)} \\
&= \frac{1}{N}\,\frac{\sin\!\big(\pi[k-\ell]\big)\,e^{j\pi(k-\ell)}}{\sin\!\big(\frac{\pi}{N}[k-\ell]\big)\,e^{j\frac{\pi(k-\ell)}{N}}}
\end{align*}
\noindent\HandoutMarkLeft{3}Hence
\begin{equation*}
(\ast) = \begin{cases} 1, & k-\ell = rN,\ r \text{ integer}, \\ 0, & \text{else} \end{cases}
\;=\; \delta\big[\HandoutMod{k-\ell}{N}\big].
\end{equation*}

\begin{notebox}[title=\HandoutNoteAddendumTitle]
The geometric-series formula is only valid when $W_N^{-(k-\ell)}\ne 1$, i.e., $k-\ell\ne rN$; in that case $\sin(\pi[k-\ell])=0$ while $\sin(\frac{\pi}{N}[k-\ell])\ne 0$, so $(\ast)=0$.
When $k-\ell = rN$, every term of the sum equals $W_N^{-rNn}=1$, so $(\ast) = \frac{1}{N}\cdot N = 1$ directly (this is also the $0/0$ limit of the Dirichlet-kernel expression).
\end{notebox}

\begin{defbox}[title={Modulo Operation $\HandoutMod{n}{N}$}]
$k = \HandoutMod{n}{N}$ if $k\in\{0,\dots,N-1\}$ and there exists an integer $\ell$ such that $n = k + \ell N$.
\end{defbox}

\begin{examplebox}[title={Example ($N=3$)}]
\begin{center}
\renewcommand{\arraystretch}{1.2}
\begin{tabular}{c|ccccccccc}
$n$ & $-2$ & $-1$ & $0$ & $1$ & $2$ & $3$ & $4$ & $5$ & $6$ \\
\hline
$\HandoutMod{n}{3}$ & $1$ & $2$ & $0$ & $1$ & $2$ & $0$ & $1$ & $2$ & $0$
\end{tabular}
\end{center}
\end{examplebox}

\noindent Therefore
\begin{equation*}
X[\ell] = \sum_{k=0}^{N-1} X[k]\,\delta\big[\HandoutMod{k-\ell}{N}\big] = X[\ell] \quad \text{for } \ell\in\{0,\dots,N-1\}. \quad \blacksquare
\end{equation*}

\begin{examplebox}[title={Example}, handout mark=4]
$g[n] = \cos\!\left(\dfrac{2\pi rn}{N}\right)$, $\;0\le n\le N-1$, with $r$ an integer, $0\le r\le N-1$.
\begin{equation*}
g[n] = \tfrac{1}{2}\left(W_N^{-rn} + W_N^{rn}\right)
\end{equation*}
\begin{align*}
\Rightarrow\ G[k] &= \frac{1}{2}\left(\sum_{n=0}^{N-1} W_N^{-(r-k)n} + \sum_{n=0}^{N-1} W_N^{-(-r-k)n}\right) \\
&= \frac{1}{2}\Big(N\,\delta\big[\HandoutMod{r-k}{N}\big] + N\,\delta\big[\HandoutMod{-r-k}{N}\big]\Big).
\end{align*}
For $0\le k\le N-1$:
\begin{equation*}
G[k] = \begin{cases} \dfrac{N}{2}, & k = r,\ \ k = N-r, \\[4pt] 0, & \text{else}, \end{cases}
\qquad \text{(unless $r=0$ or $r=N/2$)}
\end{equation*}
using $\HandoutMod{-r-k}{N} = \HandoutMod{N-r-k}{N}$. When $r=0$ or $r=N/2$, the two impulses coincide and $G[r] = N$.
\end{examplebox}

\HandoutMarkSection{5}{\HandoutTitleWithFnmark{DFT Properties}}
\footnotetext{cf.\ Mitra Tables~5.1, 5.2, and~5.3 (reproduced in the appendix of the full build).}

\subsection{Periodicity}

\begin{factbox}[title={Property}]
\begin{align*}
X[\ell N + k] &= X[k], \qquad k=0,\dots,N-1,\ \ \ell \text{ an integer}, \\
x[mN + n] &= x[n], \qquad n=0,\dots,N-1,\ \ m \text{ an integer},
\end{align*}
when the DFT / IDFT formulas are evaluated outside $0,\dots,N-1$.
\end{factbox}
Because
\begin{equation*}
W_N^N = e^{-j\frac{2\pi}{N}N} = e^{-j2\pi} = 1
\quad\Rightarrow\quad
W_N^{\ell N+k} = (1)^{\ell}\,W_N^{k} = W_N^{k}.
\end{equation*}

\subsection{Circular Shift}

Consider $x[n]$, $0\le n\le N-1$. The ordinary shift $x_1[n] = x[n-n_0]$ is \textbf{no longer defined} over $0\le n\le N-1$.
\begin{equation*}
\Rightarrow\ \text{circular shifts:}\qquad x_1[n] = x\big[\HandoutMod{n-n_0}{N}\big].
\end{equation*}

\begin{examplebox}[title={Example (Right Circular Shift)}, handout mark=6]
Let $0<n_0\le N-1$:
\begin{equation*}
x_1[n] = x\big[\HandoutMod{n-n_0}{N}\big] = \begin{cases} x[n-n_0], & n_0\le n\le N-1, \\ x[N+n-n_0], & 0\le n\le n_0-1. \end{cases}
\end{equation*}

\begin{handouttikz}[x=0.62cm, y=0.36cm]
  \foreach \n/\v in {0/1, 1/2, 2/3, 3/4} {
    \draw[thick] (\n,0) -- (\n,\v); \fill (\n,\v) circle (2pt);
    \node[above, font=\scriptsize] at (\n,\v) {$\v$};
    \node[below, font=\scriptsize] at (\n,0) {$\n$};
  }
  \draw[->] (-0.6,0) -- (4.0,0) node[right] {$n$};
  \node[font=\small] at (1.5,-2.1) {$x[n]$};
  \draw[->, thick] (4.9,2) -- (6.3,2) node[midway, above, font=\small] {$n_0=1$};
  \begin{scope}[xshift=4.4cm]
    \foreach \n/\v in {0/4, 1/1, 2/2, 3/3} {
      \draw[thick] (\n,0) -- (\n,\v); \fill (\n,\v) circle (2pt);
      \node[above, font=\scriptsize] at (\n,\v) {$\v$};
      \node[below, font=\scriptsize] at (\n,0) {$\n$};
    }
    \draw[->] (-0.6,0) -- (4.0,0) node[right] {$n$};
    \node[font=\small] at (1.5,-2.1) {$x\big[\HandoutMod{n-1}{4}\big]$};
  \end{scope}
\end{handouttikz}

\begin{equation*}
x[n] = \{\underset{\uparrow}{1},2,3,4\}, \qquad
x\big[\HandoutMod{n-1}{4}\big] = \{\underset{\uparrow}{4},1,2,3\},
\end{equation*}
\begin{equation*}
x\big[\HandoutMod{n-2}{4}\big] = \{\underset{\uparrow}{3},4,1,2\}
= x\big[\HandoutMod{n+2}{4}\big].
\end{equation*}
\end{examplebox}

\begin{factbox}[title={Property (Circular Time-Shifting)}]
Let $g[n] \longleftrightarrow G[k]$ and $\tilde g[n] = g\big[\HandoutMod{n-n_0}{N}\big]$. Then
\begin{equation*}
\tilde G[k] = W_N^{kn_0}\,G[k] = e^{-j\frac{2\pi kn_0}{N}}\,G[k].
\end{equation*}
\end{factbox}

\HandoutMarkParagraph{7}{Proof.}
\begin{align*}
\tilde G[k] &= \sum_{n=0}^{N-1} g\big[\HandoutMod{n-n_0}{N}\big]\,W_N^{kn}
= \sum_{n=0}^{n_0-1} g[n+N-n_0]\,W_N^{kn} + \sum_{n=n_0}^{N-1} g[n-n_0]\,W_N^{kn}.
\end{align*}
Let $m = n+N-n_0$ ($n = m-N+n_0$) in the first sum and $p = n-n_0$ ($n=p+n_0$) in the second:
\begin{align*}
\tilde G[k] &= \sum_{m=N-n_0}^{N-1} g[m]\,W_N^{k(m-N+n_0)} + \sum_{p=0}^{N-1-n_0} g[p]\,W_N^{k(p+n_0)} \\
&= \sum_{n=0}^{N-1} g[n]\,W_N^{k(n+n_0)} \qquad \left(W_N^{-kN}=1\right) \\
&= W_N^{kn_0}\left[\sum_{n=0}^{N-1} g[n]\,W_N^{kn}\right]
= W_N^{kn_0}\,G[k]
= e^{-j\frac{2\pi kn_0}{N}}\,G[k]. \quad \blacksquare
\end{align*}

\HandoutMarkSubsection{8}{Convolution}

\noindent Linear convolution and the DTFT:
\begin{equation*}
x[n]\circledast h[n] \ \overset{\text{DTFT}}{\longleftrightarrow}\ X(e^{j\omega})\,H(e^{j\omega}).
\end{equation*}
If $x[n]$ has length $N_1$ and $h[n]$ has length $N_2$, then $y[n] = x[n]\circledast h[n]$ has length $N_1+N_2-1$: it is possibly nonzero for $0\le n\le N_1+N_2-2$ and is zero otherwise.

\begin{defbox}[title={Circular Convolution (two length-$N$ sequences)}]
\begin{equation*}
y_c[n] = \sum_{m=0}^{N-1} x[m]\,h\big[\HandoutMod{n-m}{N}\big], \qquad 0\le n\le N-1,
\end{equation*}
\begin{equation*}
y_c[n] = x[n] \HandoutCircConv h[n] = h[n] \HandoutCircConv x[n].
\end{equation*}
\end{defbox}

\begin{factbox}[title={Convolution Theorem}]
If $y_c[n] = x[n] \HandoutCircConv h[n]$, then $Y_c[k] = X[k]\,H[k]$.
\end{factbox}

\HandoutMarkParagraph{9}{Proof.}
\begin{align*}
Y_c[k] &= \sum_{n=0}^{N-1}\left[\sum_{m=0}^{N-1} x\big[\HandoutMod{n-m}{N}\big]\,h[m]\right] W_N^{kn}
= \sum_{m=0}^{N-1} h[m]\left[\sum_{n=0}^{N-1} x\big[\HandoutMod{n-m}{N}\big]\,W_N^{kn}\right] \\
&= \left[\sum_{m=0}^{N-1} h[m]\,W_N^{km}\right] X[k] \qquad \text{(circular time-shift)} \\
&= H[k]\,X[k]. \quad \checkmark
\end{align*}

\subsection{Linear Convolution via Circular Convolution}

If $x[n]$, $0\le n\le N_1-1$, and $h[n]$, $0\le n\le N_2-1$, then
\begin{equation*}
y[n] = x[n]\circledast h[n], \qquad 0\le n\le N_1+N_2-2.
\end{equation*}

\begin{factbox}[title={Property}]
If
\begin{equation*}
\tilde x[n] = \begin{cases} x[n], & 0\le n\le N_1-1, \\ 0, & N_1\le n< N, \end{cases}
\qquad \text{where } N\ge N_1+N_2-1,
\end{equation*}
and similarly define a ``zero-padded'' length-$N$ $\tilde h[n]$, then
\begin{equation*}
\tilde x[n] \HandoutCircConv \tilde h[n] = x[n]\circledast h[n] \qquad \text{for } 0\le n\le N-1.
\end{equation*}
\end{factbox}

\begin{factbox}[title={Consequence: we can compute $\circledast$ using $\HandoutCircConv$!}, handout mark=10]
\begin{enumerate}
  \item Zero-pad both sequences to the same length $N\ge N_1+N_2-1$.
  \item Take length-$N$ DFTs of both.
  \item Multiply the DFTs together.
  \item Take the inverse DFT of the result.
\end{enumerate}
\end{factbox}

\begin{notebox}[title=\HandoutNoteAddendumTitle]
Why $N\ge N_1+N_2-1$: circular convolution equals the linear convolution \emph{aliased} with period $N$,
$\tilde x[n]\HandoutCircConv\tilde h[n] = \sum_{r} y[n+rN]$ for $0\le n\le N-1$. Since $y[n]$ occupies $N_1+N_2-1$ samples, no shifted copies overlap exactly when $N\ge N_1+N_2-1$.
With a fast algorithm (FFT) for the DFT, this is how long convolutions are computed in practice (cf.\ Mitra Ch.~11).
\end{notebox}

\subsection{Parseval's Theorem}

\begin{factbox}[title={Property}]
\begin{equation*}
\sum_{n=0}^{N-1} x^*[n]\,y[n] = \frac{1}{N}\sum_{k=0}^{N-1} X^*[k]\,Y[k]
\quad\Rightarrow\quad
\sum_{n=0}^{N-1} |x[n]|^2 = \frac{1}{N}\sum_{k=0}^{N-1} |X[k]|^2.
\end{equation*}
\end{factbox}

\subsection{Symmetry}

\begin{factbox}[title={Property}]
If $x[n]$ is real,
\begin{equation*}
X[k] = X^*[N-k] = X^*\big[\HandoutMod{-k}{N}\big], \qquad 0\le k\le N-1.
\end{equation*}
\end{factbox}
\noindent(For $k=0$, read $X^*[N]$ as $X^*[0]$ by periodicity; the modulo form handles this automatically.)

\section*{Readings}
\begin{itemize}
  \item Sanjit K. Mitra, \textit{Digital Signal Processing: A Computer-Based Approach}, 4th ed., Ch.\ 5.4, 5.5, 5.7, 11.3--11.5
  \item Monson H. Hayes, \textit{Schaum's Outline of Digital Signal Processing}, 2nd ed., Ch.\ 6, 7
\end{itemize}

\HandoutAppendixListStart
  \HandoutAppendixListItem{app:dsp-foundations}{A.\ DSP Foundations (Euler, Analytic Signals, Norms)}
  \HandoutAppendixListItem{app:mitra-tables5}{B.\ Mitra Tables 5.1, 5.2, and 5.3}
  \HandoutAppendixListItem{app:scan8}{C.\ Handwritten Lecture Note Scan (September 17, 2026)}
\HandoutAppendixListEnd

\HandoutAppendixSection{app:dsp-foundations}{DSP Foundations (Euler, Analytic Signals, Norms)}
\ifhandoutclean\else
% Course-wide reference: Euler identities, analytic signals, common norms.
% Included from ee483_lectures_clean.tex and full lecture handouts (not per-lecture clean PDFs).

\subsection*{Euler's formula and variants}

\begin{equation*}
e^{j\theta} = \cos\theta + j\sin\theta
\qquad\text{(Euler's formula).}
\end{equation*}

\noindent Identities used in GLP proofs, Fourier pairs, and Hilbert/analytic-signal material:
\begin{align*}
e^{j\theta}+e^{-j\theta} &= 2\cos\theta, &
\cos\theta &= \tfrac{1}{2}\bigl(e^{j\theta}+e^{-j\theta}\bigr), \\
e^{j\theta}-e^{-j\theta} &= 2j\sin\theta, &
\sin\theta &= \frac{1}{2j}\bigl(e^{j\theta}-e^{-j\theta}\bigr), \\
e^{j\pi/2} &= j, \qquad e^{-j\pi/2} &= -j.
\end{align*}

\subsection*{Analytic signal}

\begin{defbox}[title={Analytic signal}]
A signal whose spectrum has \textbf{no negative-frequency components} is called an \textbf{analytic signal}.

\medskip
\noindent\textbf{Continuous time.}\quad For real $x(t)$ with Fourier transform $X(j\Omega)$, form
\begin{equation*}
z(t) = x(t) + j\,\hat{x}(t), \qquad \hat{x}(t)=\mathcal{H}\{x(t)\},
\end{equation*}
where $\mathcal{H}\{\cdot\}$ is an ideal Hilbert transformer ($\mp j$ on $\Omega\gtrless0$). Then $z(t)$ is analytic: $Z(j\Omega)=0$ for $\Omega<0$ (positive frequencies retained; cf.\ Lecture~7).

\medskip
\noindent\textbf{Discrete time.}\quad Similarly $z[n]=x[n]+j\,\hat{x}[n]$ with $\hat{x}[n]=x[n]\circledast h_{\mathrm{HT}}[n]$ and ideal $H_{\mathrm{HT}}(e^{j\omega})=\mp j$ on $[0,\pi)$ / $(-\pi,0)$ gives $Z(e^{j\omega})=0$ on $-\pi\le\omega<0$.
\end{defbox}

\noindent\smallskip
\noindent\small In-phase / quadrature: $x$ is the \textbf{in-phase} and $\hat{x}$ the \textbf{quadrature} component of $z$. Often $|z(t)|$ is the \textbf{instantaneous envelope} and $\arg z(t)$ the instantaneous phase (narrowband / AM settings).

\subsection*{Common norms and energy}

\noindent\textbf{Discrete-time sequences} $x[n]$ (when the sums converge):
\begin{align*}
\|x\|_1 &= \sum_{n=-\infty}^{\infty} |x[n]|, &
\|x\|_2 &= \sqrt{\sum_{n=-\infty}^{\infty}|x[n]|^2}, &
\|x\|_\infty &= \sup_{n\in\mathbb{Z}} |x[n]|.
\end{align*}

\noindent\textbf{Continuous-time signals} $x(t)$ (when the integrals converge):
\begin{align*}
\|x\|_1 &= \int_{-\infty}^{\infty}|x(t)|\,dt, &
\|x\|_2 &= \sqrt{\int_{-\infty}^{\infty}|x(t)|^2\,dt}, &
\|x\|_\infty &= \sup_{t\in\mathbb{R}} |x(t)|.
\end{align*}

\begin{factbox}[title={Parseval (DTFT)}]
When the DTFT exists in the usual sense,
\begin{equation*}
\sum_{n=-\infty}^{\infty}|x[n]|^2
= \frac{1}{2\pi}\int_{-\pi}^{\pi}\bigl|X(e^{j\omega})\bigr|^2\,d\omega
= \|X\|_2^2 \quad \text{(energy in time = energy in frequency).}
\end{equation*}
\end{factbox}

\noindent\small\textit{Ref.:}\ J.\ O.\ Smith, ``The Analytic Signal and Hilbert Transform Filters'' (\S5.3.4), \url{https://ccrma.stanford.edu/~jos/st/Analytic_Signals_Hilbert_Transform.html}.

\clearpage
\fi

\HandoutAppendixSection{app:mitra-tables5}{Mitra Tables 5.1, 5.2, and 5.3}
\ifhandoutclean\else
% Mitra 4th ed.: Tables 5.1, 5.2, and 5.3 (DFT properties appendix).

\begin{table}[H]
\centering
\caption{Symmetry properties of the DFT of a complex sequence (Mitra Table~5.1).}
\label{tab:dft-complex-symmetry}

\begin{tabular}{cc}
\toprule
\textbf{Length-$N$ Sequence} & \textbf{$N$-point DFT} \\
\midrule
$x[n]$ & $X[k]$ \\
\midrule
$x^{*}[n]$ & $X^{*}[\HandoutMod{-k}{N}]$ \\[0.8em]
$x^{*}[\HandoutMod{-n}{N}]$ & $X^{*}[k]$ \\[0.8em]
$\operatorname{Re}\{x[n]\}$ &
$X_{\mathrm{pcs}}[k]=\dfrac{1}{2}\left\{X[\HandoutMod{k}{N}]+X^{*}[\HandoutMod{-k}{N}]\right\}$ \\[0.8em]
$j\operatorname{Im}\{x[n]\}$ &
$X_{\mathrm{pca}}[k]=\dfrac{1}{2}\left\{X[\HandoutMod{k}{N}]-X^{*}[\HandoutMod{-k}{N}]\right\}$ \\[0.8em]
$x_{\mathrm{pcs}}[n]$ & $\operatorname{Re}\{X[k]\}$ \\[0.8em]
$x_{\mathrm{pca}}[n]$ & $j\operatorname{Im}\{X[k]\}$ \\
\bottomrule
\end{tabular}

\medskip

\parbox{0.9\textwidth}{
\textit{Note:}
$x_{\mathrm{pcs}}[n]$ and $x_{\mathrm{pca}}[n]$ are the periodic
conjugate-symmetric and periodic conjugate-antisymmetric parts of $x[n]$,
respectively. Likewise, $X_{\mathrm{pcs}}[k]$ and $X_{\mathrm{pca}}[k]$ are
the periodic conjugate-symmetric and periodic conjugate-antisymmetric parts
of $X[k]$, respectively.
}
\end{table}

\begin{table}[H]
\centering
\caption{Symmetry properties of the DFT of a real sequence (Mitra Table~5.2).}
\label{tab:dft-real-symmetry}

\begin{tabular}{cc}
\toprule
\textbf{Length-$N$ Sequence} & \textbf{$N$-point DFT} \\
\midrule
$x[n]$ & $X[k]=\operatorname{Re}\{X[k]\}+j\operatorname{Im}\{X[k]\}$ \\
\midrule
$x_{\mathrm{pe}}[n]$ & $\operatorname{Re}\{X[k]\}$ \\
$x_{\mathrm{po}}[n]$ & $j\operatorname{Im}\{X[k]\}$ \\
\midrule
\multirow{5}{*}{Symmetry relations}
& $X[k]=X^{*}[\HandoutMod{-k}{N}]$ \\[0.6em]
& $\operatorname{Re}X[k]=\operatorname{Re}X[\HandoutMod{-k}{N}]$ \\[0.6em]
& $\operatorname{Im}X[k]=-\operatorname{Im}X[\HandoutMod{-k}{N}]$ \\[0.6em]
& $\lvert X[k]\rvert=\lvert X[\HandoutMod{-k}{N}]\rvert$ \\[0.6em]
& $\arg X[k]=-\arg X[\HandoutMod{-k}{N}]$ \\
\bottomrule
\end{tabular}

\medskip

\parbox{0.9\textwidth}{
\textit{Note:}
$x_{\mathrm{pe}}[n]$ and $x_{\mathrm{po}}[n]$ are the periodic even and
periodic odd parts of $x[n]$, respectively.
}
\end{table}

\begin{table}[H]
\centering
\caption{General properties of the DFT (Mitra Table~5.3).}
\label{tab:dft-theorems}

\small
\begin{tabular}{@{}>{\raggedright\arraybackslash}p{0.24\linewidth}@{}>{\raggedright\arraybackslash}p{0.34\linewidth}@{}>{\raggedright\arraybackslash}p{0.40\linewidth}@{}}
\toprule
\textbf{Type of Property} & \textbf{Length-$N$ Sequence} & \textbf{$N$-point DFT} \\
\midrule
& $g[n]$ & $G[k]$ \\
& $h[n]$ & $H[k]$ \\
\midrule
Linearity &
$\alpha g[n]+\beta h[n]$ &
$\alpha G[k]+\beta H[k]$ \\[0.6em]
Circular time-shifting &
$g[\HandoutMod{n-n_0}{N}]$ &
$W_N^{kn_0}\,G[k]$ \\[0.6em]
Circular frequency-shifting &
$W_N^{-k_0 n}\,g[n]$ &
$G[\HandoutMod{k-k_0}{N}]$ \\[0.6em]
Duality &
$G[n]$ &
$N\,g[\HandoutMod{-k}{N}]$ \\[0.6em]
$N$-point circular convolution &
$\displaystyle\sum_{m=0}^{N-1} g[m]\,h[\HandoutMod{n-m}{N}]$ &
$G[k]\,H[k]$ \\[1.2em]
Modulation &
$g[n]\,h[n]$ &
$\displaystyle\frac{1}{N}\sum_{m=0}^{N-1} G[m]\,H[\HandoutMod{k-m}{N}]$ \\[1.2em]
Parseval's relation &
\multicolumn{2}{l}{$\displaystyle\sum_{n=0}^{N-1} |x[n]|^2=\frac{1}{N}\sum_{k=0}^{N-1} |X[k]|^2$} \\
\bottomrule
\end{tabular}
\end{table}

\clearpage
\fi

\HandoutAppendixSection{app:scan8}{Handwritten Lecture Note Scan}
\HandoutIncludeHandoutScan{lecture8_0917.pdf}

\end{document}
